\documentclass[../../../include/open-logic-section]{subfiles}

\begin{document}

\olfileid{sth}{ordinals}{wo}
\olsection{సుక్రమాలు}

ప్రాథమిక భావన ఇది:

\begin{defn}
$<$ సంబంధం $A$ను
\emph{సుక్రమితం చేస్తుంది}
అంటే ఈ రెండు షరతులను
పూర్తి చేస్తుంది:
\begin{enumerate}
	\item $<$ సంపర్కితమైనది; అంటే, ప్రతి $a, b \in A$కూ $a < b$
	లేదా $a = b$ లేదా $b < a$;
	\item $A$లోని ప్రతి ఖాళీ కాని ఉపసమితికి ఒక $<$-కనిష్ఠ
	\tetoken{మూలకం}{!!{element}} ఉంటుంది; అంటే, $\emptyset \neq X \subseteq A$ అయితే $(\exists m \in
	X)(\forall z \in X)z \nless m$.
\end{enumerate}
\end{defn}

ఇప్పుడే పరిశీలించిన మూడు
ఉదాహరణలు నిజంగా
సుక్రమ సంబంధాలని
చూడడం సులభం.

\begin{prob}
\Olref[sth][ordinals][idea]{sec}
సహజ సంఖ్యలపై మూడు
ఉదాహరణ క్రమాలను
ఇచ్చింది. ప్రతిదీ
సుక్రమమని తనిఖీ
చేయండి.
\end{prob}

సుక్రమం గురించి కొన్ని
ప్రాథమికమైన, అయితే
అత్యంత ముఖ్యమైన
విషయాలు ఇవి.

\begin{prop}\ollabel{wo:strictorder}
$<$, $A$ను సుక్రమితం
చేస్తే, $A$లోని ప్రతి
ఖాళీ కాని ఉపసమితికి
ఏకైక $<$-అత్యల్ప
మూలకం ఉంటుంది; ఇంకా
$<$ స్వావర్తనరహిత,
అసమమిత, సంక్రామక
సంబంధం అవుతుంది.
\end{prop}

\begin{proof}
$X$, $A$లోని ఖాళీ కాని
ఉపసమితి అయితే, దానికి
ఒక $<$-కనిష్ఠ
\tetoken{మూలకం}{!!{element}}
$m$ ఉంటుంది; అంటే
$(\forall z \in X)z \nless m$.
$<$ సంపర్కితమైనది
కాబట్టి,
$(\forall z \in X)m \leq z$.
అందువల్ల $m$,
$X$లోని $<$-అత్యల్ప
\tetoken{మూలకం}{!!{element}}.

స్వావర్తనరాహిత్యం కోసం
$a \in A$ను స్థిరపరచండి;
$\{a\}$లోని $<$-అత్యల్ప
\tetoken{మూలకం}{!!{element}}
$a$నే, కాబట్టి
$a \nless a$.
సంక్రామకత్వం కోసం,
$a < b < c$ అయితే
$\{a, b, c\}$కు
$<$-అత్యల్ప
\tetoken{మూలకం}{!!{element}}
ఉంది కాబట్టి $a < c$.
అసమమితత్వం
స్వావర్తనరాహిత్యం,
సంక్రామకత్వం నుంచి
వస్తుంది.
\end{proof}

\begin{prop}\ollabel{propwoinduction}
$<$, $A$ను సుక్రమితం
చేస్తే, ఏ $\phi(x)$
సూత్రానికైనా:
\[
	\text{ఒకవేళ }(\forall a \in A)((\forall b < a)\phi(b) \lif 
		\phi(a))\text{ అయితే }(\forall a \in A)\phi(a).
\]
\end{prop}

\begin{proof}
విపర్యయాన్ని నిరూపిస్తాం.
$\lnot(\forall a \in
A)\phi(a)$ అనుకుందాం;
అంటే $X = \Setabs{x \in A}{\lnot\phi(x)} \neq
\emptyset$. అప్పుడు
$X$లో ఒక $<$-కనిష్ఠ
\tetoken{మూలకం}{!!{element}},
$a$, ఉంటుంది. కాబట్టి
$(\forall
b < a)\phi(b)$ కానీ
$\lnot \phi(a)$.
\end{proof}\noindent
ఈ చివరి ధర్మం సహజ
సంఖ్యలపై బలమైన ఆగమన
సూత్రాన్ని గుర్తుచేయాలి:
$(\forall n \in \omega)((\forall m < n)\phi(m)
\lif \phi(n))$ అయితే
$(\forall n \in \omega)\phi(n)$.
ఈ ధర్మం సుక్రమాన్ని
చాలా \emph{దృఢమైన}
భావనగా చేస్తుంది.\footnote{గుర్తుచేస్తున్నాం:
వేరుగా స్పష్టంగా చెప్పనంతవరకు,
అన్ని సూత్రాల్లో పరామితులు
ఉండవచ్చు.}

\end{document}
